\documentclass[11pt]{article}
\usepackage[margin=1.05in]{geometry}
\usepackage{amsmath,amssymb}
\usepackage{booktabs}
\usepackage{enumitem}
\usepackage[hidelinks]{hyperref}

\title{The Stability Predictor:\\ \large What the Video Model Is, How It Was Pretrained,
What We Train on Top of It, and the Training Plan on Our Labeled Data}
\author{Aryan Goyal}
\date{July 18, 2026 (updated July 20)}

\begin{document}
\maketitle

\begin{abstract}
\noindent
This document describes the video model at the center of the project: the stability
predictor. It explains what kind of model it is, how its backbone was pretrained,
what the final trained layer looks like, exactly how its output is used next to the
policy (and why it is never used inside policy training), the detailed plan for
training it on our labeled stamps (about 155{,}000 labeled, now all in the
training pool), and the measured results so far. Stage 1b: the pooled head reaches
AUROC 0.87 on confident held-out stamps (0.89 with LIBERO added), single-task
training does not transfer across tasks, and a single-frame DINOv2 ablation
matches the video head almost exactly, so the signal is scene configuration
rather than motion. A head trained on the closed-loop labels reaches only
AUROC 0.627, so closed-loop expansion is at best weakly predictable from the
same inputs. Stage 1c rollouts, three seeds pooled: the predictor-driven
chunk switch ties the best fixed $k$ and the best contact-oracle cell on
tool\_hang (0.800 against 0.807) while avoiding the per-step-replanning failure
mode, and stays inside the flat band on the ceiling task square. The design
principle throughout: the backbone stays frozen, the head is tiny, and the policy
is never touched.
\end{abstract}

\tableofcontents
\newpage

%=====================================================================
\section{What the model is}

The stability predictor is a video classifier with one job: look at the last 16
camera frames plus the robot's joint readings, and say whether the current moment is
open-loop stable or open-loop unstable, with a number attached (the predicted
divergence rate $\hat\lambda$).

It has two parts:

\begin{itemize}[leftmargin=2em]
  \item \textbf{A frozen backbone:} V-JEPA 2, a ViT-L video transformer
    (checkpoint \texttt{facebook/vjepa2-vitl-fpc64-256}), about 300 million
    parameters. We never update a single one of them. It converts a short video
    clip into a grid of feature tokens.
  \item \textbf{A trained head:} attentive pooling plus a two-layer MLP, a few
    million parameters. This is the only thing we train in the entire project.
    It reads the backbone's tokens and produces the stability outputs.
\end{itemize}

Why this shape: stability is a \emph{dynamics} property. What makes a moment unstable
is contact, incipient slip, and constrained motion, and these are \emph{motion} cues.
They are visible in a clip and mostly invisible in a single frame. So the backbone
must understand motion, and V-JEPA 2 is pretrained specifically to predict motion
(next section). A single-frame backbone (DINOv2) with the identical head is kept as
the ablation that tests this claim: if the video model does not beat the single-frame
model, the temporal story is wrong.

One clarification of naming, since it causes confusion: in this project the phrase
``world model'' is sometimes used loosely for this predictor. Strictly, the predictor
does not model the world; it never generates or imagines future frames. It reads
recent frames and classifies the present regime. A true action-conditioned world
model (V-JEPA 2-AC or DINO-WM) enters the project only later, on real robots, as a
replacement for the simulator in label generation (Section \ref{sec:later}).

\paragraph{Why do we need a predictor at all?} The ground-truth stability labels
are indexed by demonstration and timestep, but at execution time the robot creates a
brand new trajectory, visiting states that sit on no demonstration's timeline, so
there is nothing to look up. The regime of the current moment must be judged live.
In simulation this can be done exactly by freezing the state and running the
perturb-and-replay probe; that oracle is used in Stage 1a to test whether the
switching rule helps when given perfect information. A real robot cannot freeze time
and rewind itself to measure its own divergence, and the only stability-relevant
signal available at deployment is what the robot can see: the last second or two of
camera frames. The predictor exists purely to read the regime out of those frames.
If the oracle rule does not help, no predictor can; if it does, the predictor is
what carries the rule out of the simulator.

%=====================================================================
\section{How the backbone was pretrained}

V-JEPA 2 is self-supervised on internet-scale video (on the order of a million hours
of clips), with no labels and no pixel reconstruction. The objective is masked
latent video prediction:

\begin{enumerate}[leftmargin=2em]
  \item Take a video clip and hide large space-time blocks of it.
  \item Encode the visible parts with the encoder.
  \item A predictor network must predict the \emph{representations} of the hidden
        blocks (as produced by a slowly-updated copy of the encoder), not their
        pixels.
\end{enumerate}

Predicting representations instead of pixels matters: the model is not rewarded for
memorizing textures or colors, it is rewarded for knowing \emph{what will be where,
and how things move}. To fill in a hidden block of frames correctly in latent space,
the encoder has to carry information about object motion, contact, and continuity
over time. That is precisely the information a stability readout needs. This is why
we expect (and test) that a small head on top of these frozen features can read out
``the gripper is about to bind against the peg'' from 16 frames.

The specific checkpoint we use was pretrained at 256 by 256 resolution with 64-frame
clips (the \texttt{fpc64-256} in the name); feeding it 16-frame clips at 256
resolution is a standard supported use. We use the encoder only and discard the
pretraining predictor network.

%=====================================================================
\section{The input and what the backbone produces}

Per labeled stamp $(demo, t)$:

\begin{itemize}[leftmargin=2em]
  \item \textbf{Video input:} the last 16 frames of the main camera ending at $t$
        (padded by repeating the first frame at episode start), resized to
        $256 \times 256$. Main camera: agentview for lift, can, square, and the
        MimicGen tasks; sideview for tool\_hang; agentview for LIBERO.
  \item \textbf{Backbone output:} the encoder splits the clip into space-time
        patches (16 by 16 spatial patches, 2 frames per temporal slot, so
        $8 \times 16 \times 16$ tokens of dimension 1024). We average-pool each
        $4 \times 4$ spatial neighborhood, leaving $8 \times 4 \times 4 = 128$
        tokens of dimension 1024 per stamp. This keeps coarse spatial layout and
        full temporal structure at 1/16 the storage.
  \item \textbf{Proprioception:} end-effector position (3), orientation quaternion
        (4), gripper joint positions (2), taken at $t$: a 9-dimensional vector.
  \item \textbf{Caching:} because the backbone is frozen, these features are
        computed once and stored (fp16). All four robomimic tasks are already
        extracted (the tool\_hang set alone is 12 GB). Head training then never
        touches the encoder or the images, which is why it takes minutes per epoch
        instead of days.
\end{itemize}

An important practical note: the cached features are indexed by $(demo, t)$ and do
not depend on the label values. When we regenerated all labels with the true
measured $\sigma_u$, the features stayed valid; training joins features to the new
labels by stamp index.

%=====================================================================
\section{The final layer: the trained head}

The head has three stages, all small:

\begin{enumerate}[leftmargin=2em]
  \item \textbf{Attentive pooling.} A single learned query vector attends over the
        128 tokens (standard cross-attention with a handful of heads) and produces
        one 1024-dimensional summary vector. In plain terms: the head learns
        \emph{where and when to look} inside the clip. This is V-JEPA's own
        standard evaluation protocol (the ``attentive probe''), reused unchanged.
  \item \textbf{Proprio fusion.} The 9-dimensional proprio vector passes through a
        small MLP and is concatenated to the pooled summary.
  \item \textbf{Output MLP.} A two-layer MLP maps the concatenation to four
        outputs:
        \begin{itemize}[leftmargin=1.5em]
          \item $p_{ol}$: probability the current moment is open-loop unstable
                (a sigmoid; this is the decision output),
          \item $\hat\lambda_{ol}$: predicted open-loop divergence rate
                (a regression; calibrates the head and enables a continuous-$k$
                rule as an ablation),
          \item $p_{cl}$, $\hat\lambda_{cl}$: the same pair for the closed-loop
                channel (trained from the start, used by the full three-case
                runtime rule and the mechanism analysis).
        \end{itemize}
\end{enumerate}

If the frozen-backbone head underfits, the planned fallback is LoRA adapters on the
last encoder blocks. Nothing in the current results suggests this is needed.

%=====================================================================
\section{How the head's output is used (and where it is NOT used)}

This is the most commonly confused point, so it gets its own section.

\textbf{The predictor's output is never used in policy training.} The policy
(Diffusion Policy now, ACT in stage 2) is trained entirely on its own, with its
standard loss, finishes, and is then frozen forever. The predictor is also trained
entirely on its own, supervised by the simulator labels. The two models never see
each other during training.

They meet only at \textbf{execution time}. The policy always predicts a chunk of 16
future actions. The open question at every replan is how many of those predicted
actions to execute before replanning. That number $k$ is a queue-management knob
outside the network. The predictor sets it:

\begin{center}
\begin{tabular}{@{}lll@{}}
\toprule
Predicted regime & Executed chunk & Meaning \\
\midrule
$p_{ol}$ low (stable / not unstable) & $k = k_{\mathrm{long}}$ (8--16) & commit; plant absorbs or \\
 & & carries errors, fewer injections \\
$p_{ol}$ high, $p_{cl}$ low & $k = 1$ & react; feedback beats the \\
 & & expansion here \\
$p_{ol}$ high, $p_{cl}$ high & $k = 1$ + flag & no execution-side fix; segment \\
 & & marked for data-side repair \\
\bottomrule
\end{tabular}
\end{center}

Hysteresis (two consecutive agreeing predictions) prevents chattering at segment
boundaries. Because every compared execution rule runs the exact same frozen policy
weights, any performance difference is attributable to the switching rule alone.
That attribution is the experiment; baking the signal into policy training would
destroy it.

The one place the predictor's output may eventually influence training is stage 2,
and it is data-side, not loss-side: segments flagged by the third row (unstable
under both channels) become the target list for noise-injected demonstrations,
which is the theory's prescribed remedy for that regime.

%=====================================================================
\section{Training plan on the datasets we have}
\label{sec:plan}

\subsection{The training pool}

Supervised pairs are (cached features, labels) per stamp. Current inventory, all
labeled with the perturb-and-replay oracle at $K{=}24$, $N{=}8$, stride 2:

\begin{center}
\small
\begin{tabular}{@{}llrll@{}}
\toprule
Dataset & Stamps & Features & $\sigma_u$ & Role in training \\
\midrule
lift & 2{,}485 & cached & measured & pool \\
can & 9{,}254 & cached & measured & pool (mostly-stable balance) \\
square & 12{,}723 & cached & measured & pool + primary eval \\
tool\_hang & 45{,}633 & cached & measured & pool + primary eval \\
MimicGen stack\_d0 & 8{,}445 & to extract & interim & pool (contact diversity) \\
MimicGen square\_d0 & 12{,}874 & to extract & borrowed & pool + transfer check \\
LIBERO-Long (10 files) & $\sim$65k & to extract & interim & pool, provisional \\
policy rollouts (lift, can) & to label & to extract & measured & pool (runtime states) \\
\bottomrule
\end{tabular}
\end{center}

Three preparation steps still pending, in order of priority:
(1) label the collected lift and can policy rollouts (including the failed ones)
exactly like demos, so the predictor sees the state distribution it will face at
runtime, not only the demo tube;
(2) regenerate image observations for the MimicGen sets (their downloads are
low-dimensional) and extract features in the MimicGen environment stack;
(3) extract LIBERO features (their files ship with images) with the extractor
adapted to LIBERO's observation keys.

\subsection{Targets and losses}

\begin{itemize}[leftmargin=2em]
  \item \textbf{Decision target:} the binary label unstable-versus-rest, taken
        AFTER deadband inheritance and median filtering. The measured regime
        proportions (about 30\% unstable, near 0\% confidently stable, 70\%
        deadband) make a three-class target pointless; the runtime rule is binary
        anyway.
  \item \textbf{Decision loss:} class-weighted binary cross-entropy on $p_{ol}$.
        With a roughly 30/70 split the weighting is mild, but it is computed per
        dataset from the measured proportions rather than assumed.
  \item \textbf{Auxiliary regression:} Huber loss on $\hat\lambda_{ol}$ against the
        raw (unfiltered) $\lambda$, weight $\alpha = 0.5$. The raw value is used
        here because regression benefits from the continuous signal, including
        inside the deadband; the known tail noise in extreme $\lambda$ values is
        why this is auxiliary and Huber rather than the decision loss and L2.
  \item \textbf{Closed-loop terms:} the same BCE + Huber pair on the closed-loop
        channel, weight $\beta = 0.5$, trained only on stamps where closed-loop
        labels exist (the stride-10 subsample on robomimic). Missing-label stamps
        simply contribute no CL gradient.
\end{itemize}

\subsection{Splits and evaluation}

\begin{itemize}[leftmargin=2em]
  \item \textbf{Within-task split:} by demonstration, never by stamp (stamps 2
        steps apart are nearly identical; splitting by stamp would leak). robomimic
        tasks use the official train/valid demo masks; other suites hold out 10\%
        of demos.
  \item \textbf{Stage 1b metrics on held-out demos:} AUROC and calibration of
        $p_{ol}$ against the oracle labels; agreement of $\hat\lambda$ with oracle
        $\lambda$ (rank correlation); qualitative alignment with contact flags.
  \item \textbf{Cross-task transfer, the key generalization test:} train the head
        on Square only, evaluate label quality on Tool-Hang. This asks whether the
        model learned \emph{stability} or a task-specific visual cue. The MimicGen
        variants add a second, easier transfer axis (human square to generated
        square).
  \item \textbf{Shortcut control:} shuffle-clip test (destroy temporal order,
        expect performance to drop toward the single-frame ablation) and the
        late-in-episode control (short tasks in the pool break the
        ``late means contact'' shortcut).
  \item \textbf{The temporal claim:} DINOv2 single-frame backbone with the
        identical head and data. The gap between it and V-JEPA 2 is the measured
        value of motion information.
\end{itemize}

\subsection{Optimization details}

AdamW on the head only (backbone frozen, features precomputed): learning rate
$10^{-3}$ with cosine decay, weight decay 0.01, batch 512 stamps, up to 50 epochs
with early stopping on validation AUROC. Balanced sampling across datasets so
tool\_hang's 45k stamps do not drown lift's 2.5k. On cached features one epoch over
the full pool is minutes on one GPU; a full training run with ablations is an
afternoon, not a compute event.

\subsection{Order of experiments}

\begin{enumerate}[leftmargin=2em]
  \item Head on the four robomimic tasks (features ready today), Stage 1b metrics.
  \item Add rollout-state stamps for lift and can; check whether runtime-state
        coverage moves held-out AUROC.
  \item Cross-task transfer (Square to Tool-Hang) and the two ablations.
  \item Add MimicGen, then LIBERO features to the pool; re-measure. LIBERO enters
        as provisional until its $\sigma_u$ is measured with a LIBERO-trained
        policy.
  \item Freeze the best head and hand it to Stage 1c: the predictor-driven $k$
        switch against fixed-$k$ and policy-confidence baselines, all on the same
        frozen policy.
\end{enumerate}

%=====================================================================
\section{Stage 1b results (July 18)}
\label{sec:stage1b}

The first two trainings from the experiment list have run to completion on cached
features. Both trained the attentive head for 40 epochs, split by demonstration,
with the class-weighted BCE on confident stamps plus the Huber regression on
$\lambda$ described above.

\subsection{What was run}

Four head trainings have completed. All share the same architecture (the attentive
head of Section 5) and the same protocol: 40 epochs, batch 256, AdamW at learning
rate $10^{-3}$ with cosine decay, split by demonstration with 20 percent of demos
held out, seed 0.

\begin{description}[leftmargin=1.6em]
\item[Run 1, pooled Panda head.] Eight feature sets: the four robomimic tasks
(true measured $\sigma_u$ labels), the two rollout-state sets (lift and can
stamps collected from policy rollouts rather than demos), and the two MimicGen
sets (stack\_d0, square\_d0). 82{,}212 train / 20{,}147 validation stamps,
pooled $\delta = 0.0355$. This is the head used by Stage 1c.
\item[Run 2, transfer test.] Trained on square only (10{,}169 / 2{,}554,
$\delta = 0.0372$), then evaluated frozen on the full tool\_hang set (45{,}633
stamps) as the cross-task generalization test.
\item[Run 3, Panda plus LIBERO, superseded.] The run-1 pool plus all ten LIBERO
files (131{,}282 / 34{,}243). Its train/validation split was drawn over the
combined pool, which invalidated the comparison against run 1 (the leak note
below); its numbers are kept only for the record.
\item[Run 3b, Panda plus LIBERO, pinned validation.] Same pool as run 3
(145{,}378 train stamps, pooled $\delta = 0.0326$) but validation pinned to
run 1's exact 280 Panda demos (20{,}147 stamps), all LIBERO in training. This is
the clean version of the LIBERO experiment.
\end{description}

\subsection{As-run configuration, where it differs from the plan above}

Three details of the plan in Section 7 were simplified in the code that actually
ran. First, the BCE and Huber losses are summed with equal weight rather than the
planned $0.5$ factors; the BCE is class-weighted (positive weight from the
confident training stamps) and computed on confident stamps only, and the Huber
($\delta_{huber}=0.1$) covers all stamps including the deadband. Second, the
closed-loop output channel is not trained yet: only about 1{,}300 closed-loop
stamps exist so far (lift, can, square; tool\_hang still labeling), too few to
supervise a channel without overfitting, so the current head has two outputs
(the unstable logit and $\hat\lambda_{ol}$) and the CL channel waits for more CL
data. Third, there is no balanced dataset sampling and no early stopping; batches
are drawn uniformly, so tool\_hang and LIBERO dominate by stamp count, with the
class-weighted BCE as the only correction. A balanced-sampling rerun is the
natural follow-up if per-task validation ever shows the small tasks being
neglected. Weight decay is $10^{-4}$.

\subsection{The numbers}

\begin{center}
\small
\begin{tabular}{@{}lcccc@{}}
\toprule
Evaluation & AUROC (all) & AUROC (confident) & $\lambda$ MAE & $n$ \\
\midrule
Run 1, held-out demos, pooled     & 0.699 & \textbf{0.868} & 0.027 & 20{,}147 \\
Run 2, held-out demos, square     & 0.768 & 0.719          & 0.028 & 2{,}554 \\
Run 2, transfer to tool\_hang     & 0.562 & 0.465          & 0.072 & 45{,}633 \\
Run 3, own mixed validation       & 0.646 & 0.880          & 0.025 & 34{,}243 \\
Run 3b, own pinned validation     & 0.727 & 0.890          & 0.027 & 20{,}147 \\
\midrule
Run 1, pinned Panda val (same code path)  & 0.703 & 0.875 & 0.027 & 20{,}147 \\
Run 3b, pinned Panda val (same code path) & \textbf{0.720} & \textbf{0.893} & 0.027 & 20{,}147 \\
\bottomrule
\end{tabular}
\end{center}

The last two rows are the like-for-like LIBERO comparison: both heads scored by
one evaluation script on the identical 20{,}147 Panda stamps that neither head
trained on. Small differences against the training-time logs above them
(0.868 versus 0.875 for run 1) come from the two code paths; comparisons are only
ever made within one code path.

``All'' scores every held-out stamp against $y = (\lambda > \delta)$; ``confident''
restricts to stamps outside the deadband ($|\lambda| > \delta$), which are the only
stamps whose binary label is proven rather than ambiguous.

\subsection{What the numbers mean}

\begin{enumerate}[leftmargin=1.6em]
\item \textbf{The pooled head works where the labels are confident.} AUROC 0.868 on
confident stamps means the head separates proven-unstable from proven-stable
moments well from 16 frames plus proprioception. The lower all-stamps number
(0.699) is expected by construction: the deadband stamps it scores against are
ambiguous in the ground truth itself, so no predictor can score them cleanly.
\item \textbf{The regression channel is informative at the scale that matters.}
$\lambda$ MAE of 0.027 sits well below the labels' typical unstable magnitude
(p95 around 0.11 to 0.13 across datasets), so thresholding $\hat\lambda$ against
$\delta$ is meaningful and this is exactly how Stage 1c uses the head.
\item \textbf{Single-task training does not transfer.} Square to tool\_hang lands
at 0.562 overall and 0.465 on confident stamps, chance level or below, with the
regression error tripling. A head trained on one task learns task-specific visual
cues, not stability. The pooled head is therefore the only candidate for control,
and the practical rule going forward is that every task family we care about
contributes some labeled data to the pool; transfer is not relied on.
\item \textbf{Consequence for Stage 1c.} The frozen run-1 head, with its training
$\delta = 0.0355$ as the default switching threshold, is the head handed to the
Stage 1c predictor-driven chunk switch.
\end{enumerate}

\subsection{The LIBERO decision: in}

The pooled retrain with all ten LIBERO files added (63k stamps, proprio
harmonized to the 9-dimensional quaternion layout) was compared against run 1 on
the identical 20{,}147 Panda validation stamps, with the validation demos pinned
to run 1's exact split so the two heads saw the same training exposure. Result:
AUROC 0.720 versus 0.703 on all stamps, 0.893 versus 0.875 on confident stamps,
$\lambda$ MAE unchanged. Adding LIBERO helps the Panda tasks modestly and hurts
nothing, so LIBERO stays in the pool despite its provisional label quality.

One methodological note, recorded because it changed a number by a lot: the first
version of this comparison drew the new head's train/validation split randomly
over the combined pool, which put most of run 1's validation demos into the new
head's training set. That leaked comparison showed 0.968, a large fake
improvement. The pinned-split retrain shows the true gain, 0.018. The trainer now
supports pinning the validation demos explicitly, and every future pool-variant
comparison uses it.

\subsection{The DINOv2 ablation: motion context adds almost nothing (July 20)}
\label{sec:dino}

Run 4 asks how much of the head's signal actually comes from the 16-frame video
clip. It replaces the V-JEPA 2 clip features with features from a single
DINOv2-large frame at the query moment: one 84 to 240 pixel camera image, resized
to 224, encoded to a CLS token plus an 8 by 8 pooled patch grid (65 tokens of
dimension 1024). Stamps, proprioception, and labels are copied stamp-for-stamp
from the run-1 pool, and the split is pinned to run 1's exact validation demos
(82{,}212 train, 20{,}147 validation, same $\delta = 0.0355$), so the comparison
is clean.

\begin{center}
\small
\begin{tabular}{@{}lccc@{}}
\toprule
 & AUROC all & AUROC confident & $\lambda$ MAE \\
\midrule
Run 1: V-JEPA 2, 16 frames   & 0.699 & 0.868 & 0.0268 \\
Run 4: DINOv2, single frame  & 0.703 & 0.859 & 0.0269 \\
\bottomrule
\end{tabular}
\end{center}

The single frame matches the video clip almost exactly: the 15 extra frames of
motion context are worth at most 0.009 confident AUROC. Three consequences.
First, the stability signal the head reads is mostly scene configuration at the
query moment (what is grasped, what is near what fixture) plus proprioception,
not motion. Second, the shuffle-clip control is moot: if temporal order carried
the signal, a single frame could not tie the clip, so that control is dropped.
Third, for deployment a single-frame encoder would serve the switch nearly as
well at a fraction of the inference cost.

One caveat on scope: the ablation changes encoder and clip length together
(DINOv2 is a different pretraining), so it bounds the joint question ``does the
switch need a video encoder with motion context'' (no, on these tasks) rather
than isolating what V-JEPA 2 itself attends to.

\subsection{Run 5: a head on the closed-loop labels (July 20)}
\label{sec:clhead}

With closed-loop labeling finished on all four robomimic tasks, run 5 trains the
same head on the closed-loop expansion rate $\lambda_{\mathrm{cl}}$ instead of
the open-loop label. V-JEPA 2 features were extracted at the 3{,}969 closed-loop
stamps (151 lift, 477 can, 644 square, 2{,}697 tool\_hang, so tool\_hang is 68
percent of the pool), with the same architecture, losses, and 40 epochs on a
3{,}150 train, 819 validation split.

The label statistics change what every metric means. Closed-loop $\lambda$ is
positive at 97.6 percent of stamps, so the automatic threshold lands at
$\delta = 0.076$, essentially the median of the labels (0.080). The
classification target is therefore relative: does this stamp expand faster than
the typical stamp, not stable versus unstable. The ``confident'' subset (true
$|\lambda| > \delta$) contains only 7 negative stamps in the entire pool, so the
reported confident AUROC of 0.986 is measured against a handful of contracting
stamps and carries no weight.

The honest numbers are AUROC 0.627 over all validation stamps and a $\lambda$
MAE of 0.043 against a threshold of 0.076. Ranking closed-loop expansion from
one stamp's features is clearly harder than the open-loop task (0.699, though
run 1 had 26 times the data, so part of the gap may be data volume). This is
consistent with the mechanism: closed-loop expansion measures how the policy
reacts to a perturbation, and that reaction depends on the policy, not only on
the scene the encoder can see. Some signal transfers, since 0.627 is well above
chance on 819 stamps, but nothing near a deployable switch. Nothing in the
pipeline currently depends on this head; it exists to answer whether the
closed-loop quantity is predictable from the same inputs, and the answer so far
is only weakly.

%=====================================================================
\section{Stage 1c rollout results (July 18 to 20)}
\label{sec:stage1c}

Stage 1c puts the frozen run-1 head in the control loop: at every replan, the last
16 frames and the current proprioception go through the encoder and head, and the
predicted $\hat\lambda$ picks the executed chunk length ($k{=}16$ if
$\hat\lambda \leq \delta$, $k{=}4$ if $\hat\lambda > \delta$, with
$\delta = 0.0355$ from training). Inference runs in a separate process with the
newer transformers stack, about 100 ms per replan. Both tasks, 50 episodes,
seed 0:

\begin{center}
\small
\begin{tabular}{@{}lcccc@{}}
\toprule
 & Predictor (16,4) & Best oracle cell & Best fixed $k$ & Fixed $k{=}1$ \\
\midrule
square     & 0.76 & 0.88 & 0.86 & 0.76 \\
tool\_hang & \textbf{0.82} & 0.78 & 0.88 & 0.62 \\
\bottomrule
\end{tabular}
\end{center}

On square everything is inside the flat band and no conclusion follows, as
expected for the ceiling task. Tool\_hang is the informative row: the predictor
cell beats every contact-oracle cell (0.82 against 0.56 to 0.78) and sits within
one-seed noise of the best fixed $k$ (0.88, gap 0.06 against a resolvable 0.16).
Its successful episodes are also the fastest measured on the task (423 steps
against 436 to 480 for all other cells).

The diagnostics say why. The predictor calls ``unstable'' on 71.6 percent of its
decisions, well below the oracle's 84 to 88 percent contact rate, and its
per-episode mean executed $k$ is 11.4. It is not reproducing the binary contact
signal; it is discriminating within contact, committing long through contact
stretches it judges benign. That is precisely the continuous-signal advantage the
Stage 1a analysis said a binary oracle could not express, now visible in a
rollout. What Stage 1c has not yet shown is the predictor beating the best fixed
$k$; on these short contact-dominated tasks that may not be possible, and the
long-horizon tasks are where that question gets a fair test. That test is now
Stage 2 (four MimicGen tasks, policies training as of July 20); its design and
status are in the Stage 2 companion document (\texttt{stage2\_doc\_latex}).

\subsection{Seed repeats and the live probe (July 19 to 20)}

The predictor cell was repeated at rollout seeds 1 and 2 on both tasks, and the
live-probe oracle cell ran once at seed 0. Pooled over three seeds
($n{=}150$, one standard error about 0.033), with the pooled Stage 1a numbers as
reference:

\begin{center}
\small
\begin{tabular}{@{}lccccc@{}}
\toprule
 & Seeds (0,1,2) & Pooled & Best fixed & Best oracle & Fixed $k{=}1$ \\
\midrule
square     & 0.76 / 0.68 / 0.86 & 0.767 & 0.833 & 0.873 & 0.793 \\
tool\_hang & 0.82 / 0.84 / 0.74 & \textbf{0.800} & 0.807 & 0.807 & 0.667 \\
\bottomrule
\end{tabular}
\end{center}

Tool\_hang holds up under pooling: the predictor at 0.800 statistically ties both
the best fixed $k$ (0.807) and the best contact-oracle cell (0.807), while
executing a mean $k$ near 11 and avoiding the 0.60 to 0.67 trap of per-step
replanning in contact. On square the predictor pools to 0.767, at the bottom of
the flat band and about 1.4 standard errors below the best fixed cell; on the
ceiling task the switch adds nothing, which is the expected null.

The live-probe cell (perturb-and-replay FTLE computed in the simulator at every
replan, probe $K{=}12$, $N{=}6$, $\sigma_u{=}0.05$, threshold $\delta{=}0.05$)
came out 0.84 on square and 0.76 on tool\_hang with mean executed $k$ of 15.8 and
15.7. At that threshold the probe almost never fired, so the cell degenerates to
fixed $k{=}16$ with probe overhead, and it is uninformative as an upper bound.
Its probe protocol ($K{=}12$) also does not share a $\delta$ scale with the
$K{=}24$ labeling protocol. Rerunning at a lower threshold would cost another
day of GPU time for a bound the predictor cell already brackets in practice, so
it is recorded and de-prioritized.

%=====================================================================
\section{The later real-robot role of a true world model}
\label{sec:later}

On a real robot there is no simulator to reset and replay, so the label source
changes: an action-conditioned world model (V-JEPA 2-AC, which adds action inputs to
the same encoder family, or DINO-WM) imagines the paired rollouts in latent space.
Same actions, perturbed first action, slope of the log latent divergence between the
imagined futures: the same label, produced by imagination instead of simulation. The
predictor's architecture, inputs, and losses do not change; only the oracle that
writes the labels does. Hardware perturb-and-replay on a subsample is the fallback.
This is deferred work and nothing in the current pipeline depends on it.

%=====================================================================
\section{Summary}

The model is a frozen 300M-parameter video transformer with a few-million-parameter
trained head. The backbone learned motion by predicting masked video in latent space
on internet-scale data; we never touch its weights. The head learns to read
stability out of 16 frames plus proprioception, supervised by the labeled stamp
pool (LIBERO included after the pinned-split comparison showed it helps). Stage
1b is now measured: the pooled head reaches AUROC 0.87 on confident held-out
stamps, single-task training does not transfer (0.56 across tasks), and a
single-frame DINOv2 head matches the video head almost exactly, so the signal is
scene configuration rather than motion. A head trained on the closed-loop labels
reaches only 0.627, consistent with closed-loop expansion being a property of
the policy's reaction rather than the visible scene. In rollouts (Stage 1c, three seeds
pooled) the predictor switch ties the best fixed and best oracle cells on
tool\_hang at 0.800 while replanning adaptively, and stays inside the flat band
on the ceiling task square.
Its output never enters policy training; it gates, at execution time, how many
actions of the frozen policy's predicted chunk are executed before replanning.
Training the head is cheap by construction: every expensive thing (backbone
pretraining, feature extraction, label generation) happens once, offline, and is
already done.

\end{document}
